\section*{Capítulo \ref{ch:b1:layered-earth} --- A Terra estratificada}

\begin{solution}{exo:b1:layered-earth:1}
Os fatores de volume são os cubos das razões dos raios: núcleo
$(3480/6371)^3 = 0.1630$, isto é, \qty{16.3}{\percent}; \omterm{def:b1:layered-earth:layers}{núcleo interno}
$(1221.5/6371)^3 = 0.00705$, \qty{0.70}{\percent}.
\end{solution}

\begin{solution}{exo:b1:layered-earth:2}
Acima: $\Phi = 13.7166^2 - \tfrac43\times7.26466^2 = 188.145 - 70.367 =
\qty{117.78}{km^2.s^{-2}}$, $K = \rho\Phi = \qty{5566.45}{kg.m^{-3}}\times
\qty{117.78e6}{m^2.s^{-2}} = \qty{655.6}{GPa}$. Abaixo: $\Phi = 8.06482^2 =
\qty{65.04}{km^2.s^{-2}}$, $K = 9903.49\times\num{65.04e6} = \qty{644.1}{GPa}$.
O módulo de compressão quase não muda (\qty{-1.7}{\percent}). O que muda é o
módulo de cisalhamento, $\mu = \rho v_S^2 = \qty{293.8}{GPa}$ acima e $0$ abaixo (um
líquido), e a densidade, que sobe \qty{78}{\percent}: daí a queda de
$v_P$ apesar de um $K$ quase igual.
\end{solution}

\begin{solution}{exo:b1:layered-earth:3}
$\sin i_c = 6.2/8.1 = 0.7654$, $i_c = \ang{49.95}$. Distância crítica
$2\times30\times\tan\ang{49.95} = 60\times1.1906 = \qty{71.4}{km}$. \omterm{def:b1:layered-earth:head-wave}{Distância de
cruzamento} $60\sqrt{14.3/1.9} = 60\times2.743 = \qty{164.6}{km}$.
\end{solution}

\begin{solution}{exo:b1:layered-earth:4}
$g = Gm/r^2 = \num{6.674e-11}\times\num{1.94e24}/(\num{3.48e6})^2 =
\qty{10.69}{m.s^{-2}}$, \qty{9}{\percent} a mais que na superfície. Descendo
até o núcleo, o raio ao quadrado cai para $(3480/6371)^2 = 0.298$ de seu valor na
superfície, enquanto a massa atrativa cai apenas para $1.94/5.972 = 0.325$ do total,
porque o núcleo contém um terço da massa: $0.325/0.298 = 1.09$.
\end{solution}

\begin{solution}{exo:b1:layered-earth:5}
Perto da fonte $v_1 = 60/10.0 = \qty{6.0}{km.s^{-1}}$. Longe
$v_2 = (300 - 200)/(45.8 - 33.3) = 100/12.5 = \qty{8.0}{km.s^{-1}}$; o
tempo de interseção é $t_i = 33.3 - 200/8.0 = \qty{8.3}{s}$. Então $\sin i_c =
0.75$, $\cos i_c = 0.6614$, e $h = t_i v_1/(2\cos i_c) = 8.3\times6.0/1.3229
= \qty{37.6}{km}$. Confere: $x_c = 2\times37.6\times\sqrt{14/2} = \qty{199}{km}$,
em que as duas ondas chegam em $199/6.0 = 199/8.0 + 8.3 = \qty{33.2}{s}$.
\end{solution}

\begin{solution}{exo:b1:layered-earth:6}
Para $\rho$ uniforme, $I = \tfrac{8\pi}{3}\rho\int_0^R r^4\,\mathrm dr =
\tfrac{8\pi}{15}\rho R^5$ e $M = \tfrac43\pi\rho R^3$, logo
$I/MR^2 = (8/15)/(4/3) = 0.4$. Duas camadas ($r_c = \qty{3.48e6}{m}$,
$R = \qty{6.371e6}{m}$): $M = \tfrac43\pi\bigl(11000\,r_c^3 + 4450\,(R^3 -
r_c^3)\bigr) = \qty{5.977e24}{kg}$, \qty{0.08}{\percent} acima da massa
medida; $I = \tfrac{8\pi}{15}\bigl(11000\,r_c^5 + 4450\,(R^5 - r_c^5)\bigr)$ dá
$I/MR^2 = 0.3457$. A massa está certa, mas o fator fica \qty{4.5}{\percent}
alto demais: o \omterm{def:b1:layered-earth:layers}{manto} real não é uniforme, sua densidade cresce com a profundidade, o que
desloca massa para dentro.
\end{solution}

\begin{solution}{exo:b1:layered-earth:7}
$\Phi = 12.16^2 - \tfrac43\times6.673^2 = 147.87 - 59.37 =
\qty{88.49}{km^2.s^{-2}}$. Com a profundidade $z$,
$\mathrm d\rho/\mathrm dz = \rho g/\Phi = 4860\times9.9/\num{88.49e6} =
\qty{5.44e-4}{kg.m^{-4}}$, isto é, \qty{0.544}{kg.m^{-3}} por quilômetro, ou
\qty{0.54}{g.cm^{-3}} por \qty{1000}{km}. O PREM sobe $5.357 - 4.844 =
0.513$ nesses \qty{1000}{km}, \qty{6}{\percent} a menos: o gradiente
calculado no topo é aplicado a todo o intervalo, enquanto $\Phi$ cresce com
a profundidade e o gradiente diminui.
\end{solution}

\begin{solution}{exo:b1:layered-earth:8}
$P \approx \bar\rho g z = 3570\times9.9\times\num{670e3} = \qty{23.7}{GPa}$
(PREM: \qty{23.9}{GPa}). É a pressão, cerca de \qty{24}{GPa}, na qual
a ringwoodita se decompõe em bridgmanita e ferropericlásio: a
descontinuidade de \qty{660}{km}.
\end{solution}

\begin{solution}{exo:b1:layered-earth:9}
(a)~Ferro líquido: $10.90\times1.05 = 11.4$ a $10.90\times1.10 =
\qty{12.0}{g.cm^{-3}}$. (b)~$12.764/12.166 = 1.049$: um salto de
\qty{4.9}{\percent} no PREM, um pouco mais que os \qty{4.5}{\percent} de
estudos recentes. (c)~A solidificação sozinha tornaria o sólido apenas um pouco mais denso
que o líquido; o salto maior significa que o sólido contém menos
elementos leves que o líquido: eles preferem o líquido, e a solidificação na
fronteira do \omterm{def:b1:layered-earth:layers}{núcleo interno} deixa a maior parte deles para trás.
\end{solution}

\begin{solution}{exo:b1:layered-earth:10}
Seja $x_* = 2h\tan i_c$. Então $\sqrt{x_*^2 + 4h^2} = 2h/\cos i_c$ e
$t_r(x_*) = 2h/(v_1\cos i_c)$. A \omterm{def:b1:layered-earth:head-wave}{onda de cabeça} dá
$t(x_*) = 2h\tan i_c/v_2 + 2h\cos i_c/v_1 = \frac{2h}{v_1}\bigl(\sin^2 i_c/\cos
i_c + \cos i_c\bigr) = 2h/(v_1\cos i_c)$, usando $\sin i_c = v_1/v_2$: o mesmo.
A inclinação da hipérbole é $t_r'(x) = x/(v_1\sqrt{x^2 + 4h^2})$; em $x_*$ ela é
$\sin i_c/v_1 = 1/v_2$, a inclinação da reta da \omterm{def:b1:layered-earth:head-wave}{onda de cabeça}. Como
$t_r'' > 0$, a hipérbole é convexa e fica acima de cada uma de suas tangentes: para
$x > x_*$, $t_r > t_{\text{cabeça}}$. Na distância crítica a reflexão
atinge a fronteira exatamente em $i_c$ e é o nascimento da \omterm{def:b1:layered-earth:head-wave}{onda de cabeça}; além,
a \omterm{def:b1:layered-earth:head-wave}{onda de cabeça}, que ganha por correr no meio rápido, chega primeiro.
\end{solution}

\begin{solution}{exo:b1:layered-earth:11}
$m(r) = \tfrac43\pi\rho r^3$, logo $g = Gm/r^2 = \tfrac43\pi G\rho r$. Então
$P(r) = \int_r^R \rho\,\tfrac43\pi G\rho s\,\mathrm ds = \tfrac23\pi G\rho^2(R^2
- r^2)$. No centro $P = \tfrac23\pi\times\num{6.674e-11}\times5513^2\times
(\num{6.371e6})^2 = \qty{172}{GPa}$, \qty{47}{\percent} do PREM
\qty{363.85}{GPa}: concentrar a massa em direção ao centro quase dobra
a pressão central.
\end{solution}

\begin{solution}{exo:b1:layered-earth:12}
Com $z$ a profundidade, $\mathrm d\rho/\mathrm dz = \rho g/\Phi$; se $g/\Phi$ é uma
constante $a$, $\mathrm d(\ln\rho)/\mathrm dz = a$ e $\rho = \rho_0
e^{a(z - z_0)}$. Aqui $a(z - z_0) = 7.9\times\num{2258.5e3}/\num{87e6} =
0.2051$ e $\rho = 9.90\times e^{0.2051} = 9.90\times1.2276 =
\qty{12.15}{g.cm^{-3}}$, \qty{0.2}{\percent} abaixo do $12.17$ do PREM: todo o
aumento de densidade através do \omterm{def:b1:layered-earth:layers}{núcleo externo} é autocompressão.
\end{solution}

\begin{solution}{pb:b1:layered-earth:1}
\textbf{1.} $V = \tfrac43\pi(\num{6.371e6})^3 = \qty{1.0832e21}{m^3}$.

\textbf{2.} $\tfrac43\pi(r_2^3 - r_1^3)$: \omterm{def:b1:layered-earth:layers}{crosta}, \omterm{def:b1:layered-earth:layers}{manto} superior e zona de transição
\qty{3.0706e20}{m^3}; \omterm{def:b1:layered-earth:layers}{manto} inferior \qty{5.9961e20}{m^3}; \omterm{def:b1:layered-earth:layers}{núcleo externo}
\qty{1.6890e20}{m^3}; \omterm{def:b1:layered-earth:layers}{núcleo interno} \qty{7.634e18}{m^3}.

\textbf{3.} \qty{28.35}{\percent}, \qty{55.36}{\percent}, \qty{15.59}{\percent}
e \qty{0.70}{\percent}: o \omterm{def:b1:layered-earth:layers}{manto} inferior é a maior casca, mais da metade
do volume da Terra.

\textbf{4.} $15.59 + 0.70 = \qty{16.3}{\percent}$.

\textbf{5.} \omterm{def:b1:layered-earth:layers}{Crosta}, \omterm{def:b1:layered-earth:layers}{manto} superior e zona de transição $\num{3.0706e20}\times3570 =
\qty{1.0962e24}{kg}$; \omterm{def:b1:layered-earth:layers}{manto} inferior \qty{2.9381e24}{kg}; \omterm{def:b1:layered-earth:layers}{núcleo externo}
\qty{1.8410e24}{kg}; \omterm{def:b1:layered-earth:layers}{núcleo interno} \qty{9.841e22}{kg}.

\textbf{6.} \qty{5.9737e24}{kg}, \qty{0.03}{\percent} acima da massa
\qty{5.972e24}{kg}.

\textbf{7.} $\num{5.9737e24}/\num{1.0832e21} = \qty{5515}{kg.m^{-3}}$.

\textbf{8.} $\num{9.841e22}/\num{5.9737e24} = \qty{1.65}{\percent}$.

\textbf{9.} $(1.0962 + 2.9381)/5.9737 = \qty{67.5}{\percent}$.

\textbf{10.} O núcleo tem \qty{16.3}{\percent} do volume e
\qty{32.5}{\percent} da massa, o dobro: sua densidade média,
\qty{10986}{kg.m^{-3}}, é cerca de $2.5$ vezes a do \omterm{def:b1:layered-earth:layers}{manto}.

\textbf{11.} Uma casca de densidade uniforme dá
$I = \tfrac{8\pi}{3}\rho\int_{r_1}^{r_2} r^4\,\mathrm dr = \tfrac{8\pi}{15}\rho
(r_2^5 - r_1^5)$.

\textbf{12.} \omterm{def:b1:layered-earth:layers}{Crosta}, \omterm{def:b1:layered-earth:layers}{manto} superior e zona de transição \qty{2.676e37}{kg.m^2}; \omterm{def:b1:layered-earth:layers}{manto} inferior
\qty{4.525e37}{kg.m^2}; \omterm{def:b1:layered-earth:layers}{núcleo externo} \qty{9.27e36}{kg.m^2}; \omterm{def:b1:layered-earth:layers}{núcleo interno}
\qty{5.87e34}{kg.m^2}.

\textbf{13.} O \omterm{def:b1:layered-earth:layers}{manto} inferior: tem metade da massa, em grandes raios. O \omterm{def:b1:layered-earth:layers}{núcleo
interno} tem \qty{1.6}{\percent} da massa, toda ela perto do eixo:
sua parcela é \qty{0.07}{\percent}.

\textbf{14.} $I = \qty{8.1345e37}{kg.m^2}$ e $I/MR^2 = \num{8.1345e37}/
(\num{5.9737e24}\times(\num{6.371e6})^2) = 0.3355$.

\textbf{15.} \qty{1.4}{\percent} acima de $0.3308$, e bem abaixo de $0.4$: o
modelo tem a concentração de massa em direção ao centro aproximadamente certa.

\textbf{16.} Sim. Dentro de cada casca a densidade real é maior embaixo do que em cima;
uma casca homogênea da mesma massa põe massa demais perto do seu raio externo, onde ela
conta com $r^2$, e assim superestima $I$.

\textbf{17.} $g = GM/R^2 = \num{6.674e-11}\times\num{5.9737e24}/
(\num{6.371e6})^2 = \qty{9.822}{m.s^{-2}}$.

\textbf{18.} A massa do núcleo é $\num{1.8410e24} + \num{9.841e22} =
\qty{1.9394e24}{kg}$, e $g = G m/r_c^2 = \qty{10.69}{m.s^{-2}}$.

\textbf{19.} $(\num{1.0962e24} + \num{2.9381e24})/(\num{3.0706e20} +
\num{5.9961e20}) = \num{4.0343e24}/\num{9.0667e20} = \qty{4450}{kg.m^{-3}}$.

\textbf{20.} $P \approx 4450\times10\times\num{2.891e6} = \qty{128.6}{GPa}$,
\qty{5}{\percent} abaixo dos \qty{135.75}{GPa} do PREM.

\textbf{21.} $\rho_c = \num{1.9394e24}/\num{1.7653e20} = \qty{10986}{kg.m^{-3}}$,
e $\Delta P = \tfrac23\pi\times\num{6.674e-11}\times10986^2\times
(\num{3.48e6})^2 = \qty{204.3}{GPa}$.

\textbf{22.} $128.6 + 204.3 = \qty{332.9}{GPa}$, \qty{8.5}{\percent} abaixo dos
\qty{363.85}{GPa} do PREM. Ambas as estimativas ignoram o aumento da densidade com a
profundidade dentro de cada região, que coloca mais massa mais fundo, mantém a gravidade
mais alta na descida e eleva $\rho g$, o gradiente de pressão.

\textbf{23.} A massa total (questão 6), o \omterm{def:b1:comparative-planetology:moi}{fator de momento de inércia}
(questão 15) e a gravidade na superfície (questão 17, contra a média medida
\qty{9.82}{m.s^{-2}}) são conferidos com medidas da Terra; a gravidade no núcleo e as
pressões só podem ser comparadas com o PREM, ele próprio
um modelo.

\textbf{24.} O núcleo tem \textbf{\qty[text-series-to-math=true]{32.5}{\percent}}
da massa da Terra.
\end{solution}
